\documentclass[11pt,a4paper]{article}

\usepackage[margin=1in]{geometry}
\usepackage{amsmath,amssymb,amsthm}
\usepackage{enumitem}
\usepackage{hyperref}
\usepackage{xcolor}
\usepackage{algorithm}
\usepackage{algpseudocode}
\usepackage{microtype}

\hypersetup{
  colorlinks=true,
  linkcolor=blue!50!black,
  urlcolor=blue!60!black,
  citecolor=blue!50!black,
}

\newtheorem{theorem}{Theorem}
\newtheorem{lemma}[theorem]{Lemma}
\newtheorem{proposition}[theorem]{Proposition}
\newtheorem{corollary}[theorem]{Corollary}
\theoremstyle{definition}
\newtheorem{definition}{Definition}
\theoremstyle{remark}
\newtheorem{remark}{Remark}

\newcommand{\name}{\mathsf{name}}
\newcommand{\rec}{\mathsf{rec}}
\newcommand{\KEM}{\mathsf{KEM}}
\newcommand{\Sig}{\mathsf{Sig}}
\newcommand{\Adv}{\mathsf{Adv}}
\newcommand{\negl}{\mathsf{negl}}
\newcommand{\KGen}{\mathsf{KGen}}
\newcommand{\Encap}{\mathsf{Encap}}
\newcommand{\Decap}{\mathsf{Decap}}
\newcommand{\Sign}{\mathsf{Sign}}
\newcommand{\Verify}{\mathsf{Verify}}
\newcommand{\pk}{\mathit{pk}}
\newcommand{\sk}{\mathit{sk}}
\newcommand{\RNS}{\textsc{ZAP-RNS}}
\newcommand{\bit}{\{0,1\}}

\title{\RNS: Formal Cryptographic Analysis}
\author{ZAP Protocol Authors\\\texttt{zap-proto.io}}
\date{2026}

\begin{document}
\maketitle

\begin{abstract}
This document accompanies the \RNS{} design paper. We give a formal
security model for the registration--lookup--connection pipeline,
present three theorems and one supporting lemma, and prove explicit
concrete bounds in terms of the underlying primitives' advantages. The
core results are: (1) name-spoofing reduces to signature forgery plus
KEM key recovery plus root compromise; (2) connection binding reduces
to KEM IND-CCA2; and (3) Sybil registration cost is bounded below by
attestation-bypass cost.
\end{abstract}

\section{Primitives and Notation}

We work in the standard concrete-security setting. Let $\lambda$ be the
security parameter.

\paragraph{KEM.} A key encapsulation mechanism is a triple
$\KEM = (\KEM.\KGen, \KEM.\Encap, \KEM.\Decap)$:
\begin{itemize}[leftmargin=*]
\item $\KEM.\KGen(1^\lambda) \to (\pk, \sk)$.
\item $\KEM.\Encap(\pk) \to (c, K)$ with $K \in \bit^\lambda$.
\item $\KEM.\Decap(\sk, c) \to K$ or $\bot$.
\end{itemize}
We use IND-CCA2 security~\cite{shoup2001} and key-recovery hardness:

\begin{definition}[KEM key-recovery]
$\Adv^{\mathrm{kr}}_{\KEM}(\mathcal{A}) = \Pr\bigl[(\pk, \sk) \gets
\KEM.\KGen(1^\lambda);\ \sk' \gets \mathcal{A}^{\Decap(\sk, \cdot)}(\pk):
\sk' = \sk \bigr]$.
\end{definition}

\begin{definition}[KEM IND-CCA2]
$\Adv^{\mathrm{ind\text{-}cca2}}_{\KEM}(\mathcal{A}) =
\bigl|\Pr[\mathbf{Exp}^{\mathrm{ind\text{-}cca2}\text{-}1}_{\KEM}(\mathcal{A}) = 1]
- \Pr[\mathbf{Exp}^{\mathrm{ind\text{-}cca2}\text{-}0}_{\KEM}(\mathcal{A}) = 1]\bigr|$,
where in experiment $b$ the challenger gives $\mathcal{A}$ a challenge
ciphertext $c^*$ and either the real shared secret ($b=1$) or a
uniform random string ($b=0$); $\mathcal{A}$ has a decapsulation oracle
on all ciphertexts $\ne c^*$.
\end{definition}

\paragraph{Signature.} A signature scheme $\Sig = (\Sig.\KGen,
\Sig.\Sign, \Sig.\Verify)$ with EUF-CMA advantage
$\Adv^{\mathrm{euf\text{-}cma}}_{\Sig}(\mathcal{A})$, defined as the
probability of producing a forgery $(m^*, \sigma^*)$ on a message not
queried to the signing oracle.

\paragraph{Hash.} $H : \bit^* \to \bit^{2\lambda}$ is modelled as a
random oracle.

\section{System Model}

A trust domain has three roles: a single \emph{root authority} holding
$(\pk_{\mathrm{root}}, \sk_{\mathrm{root}})$ for $\Sig$; a set of
\emph{services}, each potentially registering a name; and a set of
\emph{clients} performing lookups. The root key is provisioned out of
band into every client (e.g., as part of the workload image). Records
are as in Definition~1 of the design paper.

\paragraph{Registration oracle $\mathcal{O}_{\mathrm{reg}}$.} On input
$(N, \mathit{evidence})$ from an honest service the oracle generates
fresh $(\pk^\KEM, \sk^\KEM)$ and $(\pk^\Sig, \sk^\Sig)$, builds a
record $r$, signs it with $\sk^\Sig$ to produce
$\sigma_{\mathrm{owner}}$, attaches $\sigma_{\mathrm{auth}}$ from
$\sk_{\mathrm{root}}$, stores $\rec$ in a database $\mathcal{D}$, and
returns $(\rec, \pk^\KEM, \pk^\Sig)$ to the caller. Private keys are
not exposed.

\paragraph{Lookup oracle $\mathcal{O}_{\mathrm{lkp}}$.} On input $N$
returns the latest record stored for $N$ in $\mathcal{D}$, or $\bot$.

\paragraph{Adversary capability.} The adversary $\mathcal{A}$ is
polynomial-time. It may call $\mathcal{O}_{\mathrm{reg}}$ on names of
its choice, including names it claims to ``own'', and call
$\mathcal{O}_{\mathrm{lkp}}$ on any name. It may also corrupt at most a
threshold $\tau$ of services (revealing their private keys), where
$\tau$ is a parameter of the model. We do \emph{not} grant
$\mathcal{A}$ access to $\sk_{\mathrm{root}}$; root compromise is
treated as a separate event $E_{\mathrm{root}}$ with prior
probability~$p_{\mathrm{root}}$ controlled by operational policy.

\section{Theorem~1: Name-Spoofing Reduction}

\begin{definition}[Spoofing experiment]
$\mathbf{Exp}^{\mathrm{spoof}}_{\RNS}(\mathcal{A})$:
\begin{enumerate}[leftmargin=*,label=\arabic*.]
\item Challenger runs $(\pk_{\mathrm{root}}, \sk_{\mathrm{root}}) \gets
\Sig.\KGen(1^\lambda)$. $\pk_{\mathrm{root}}$ is given to $\mathcal{A}$.
\item $\mathcal{A}$ has access to $\mathcal{O}_{\mathrm{reg}}$ (which uses
$\sk_{\mathrm{root}}$) and $\mathcal{O}_{\mathrm{lkp}}$.
\item $\mathcal{A}$ outputs $(N^*, \rec^*)$.
\item Win condition: $\rec^*$ verifies in full ($\sigma_{\mathrm{owner}}$
under the embedded $\pk^{\Sig,*}$, $\sigma_{\mathrm{auth}}$ under
$\pk_{\mathrm{root}}$), and \emph{either} $N^*$ has no record in
$\mathcal{D}$ \emph{or} $\rec^*$.$\pk^\KEM \ne$
$\mathcal{D}[N^*].\pk^\KEM$.
\end{enumerate}
\end{definition}

\begin{theorem}[Name-spoofing bound]\label{thm:spoof}
For every polynomial-time adversary $\mathcal{A}$ in the random oracle
model,
\[
\Adv^{\mathrm{spoof}}_{\RNS}(\mathcal{A}) \;\le\;
\Adv^{\mathrm{euf\text{-}cma}}_{\Sig}(\mathcal{B}_1) \;+\;
q_R \cdot \Adv^{\mathrm{euf\text{-}cma}}_{\Sig}(\mathcal{B}_2) \;+\;
p_{\mathrm{root}} \;+\; \frac{q_H}{2^{2\lambda}},
\]
where $q_R$ is the number of registration-oracle calls, $q_H$ is the
number of random-oracle queries, and $\mathcal{B}_1, \mathcal{B}_2$ are
EUF-CMA adversaries built from $\mathcal{A}$ with comparable running
time.
\end{theorem}

\begin{proof}
We define a sequence of games $G_0, G_1, G_2, G_3$ and bound the
distance between consecutive games.

\paragraph{Game $G_0$.} The original spoofing experiment. Define $S_0
= \Pr[\mathcal{A} \text{ wins in } G_0]$.

\paragraph{Game $G_1$.} Identical to $G_0$ except the challenger
aborts if the event $E_{\mathrm{root}}$ occurs (root key compromise).
Then $|S_0 - S_1| \le p_{\mathrm{root}}$.

\paragraph{Game $G_2$.} Identical to $G_1$ except the challenger
aborts if any random-oracle query produces a collision: that is, if
$H(x) = H(y)$ for $x \ne y$ ever occurs in queries made by
$\mathcal{A}$ or by the oracles. By the birthday bound $|S_1 - S_2|
\le q_H^2 / 2^{2\lambda+1} \le q_H / 2^{2\lambda}$ for $q_H$ small
relative to $2^\lambda$.

\paragraph{Case split on $\mathcal{A}$'s win.} Let $\rec^*$ have
embedded public keys $(\pk^{\KEM,*}, \pk^{\Sig,*})$ and signatures
$(\sigma^*_{\mathrm{owner}}, \sigma^*_{\mathrm{auth}})$. We split the
event ``$\mathcal{A}$ wins'' into two sub-events:

\textbf{Case A.} $\sigma^*_{\mathrm{auth}}$ is on a message
$M^* = H(\rec^* \setminus \sigma^*_{\mathrm{auth}})$ that
$\mathcal{O}_{\mathrm{reg}}$ never signed. Because of $G_2$ collision
abort, $M^*$ uniquely determines a record-content. The adversary then
holds a forgery on $\sk_{\mathrm{root}}$.

\textbf{Case B.} $\sigma^*_{\mathrm{auth}}$ is on a message
$M^*$ that $\mathcal{O}_{\mathrm{reg}}$ \emph{did} sign. Then the
record-content matches one of the at most $q_R$ records previously
counter-signed. Call that record $\rec_i$. The win condition rules
out $\rec^* = \rec_i$, but in $G_2$ collisions are excluded, so the
content under $H$ matches exactly. The win condition then forces
either (i) $N^* \ne \rec_i.N$ or (ii) $\pk^{\KEM,*} \ne \rec_i.\pk^\KEM$
or (iii) $\pk^{\Sig,*} \ne \rec_i.\pk^\Sig$. All three contradict the
content equality. So Case B is impossible in $G_2$.

\paragraph{Game $G_3$.} Reduction $\mathcal{B}_1$ to EUF-CMA on
$\Sig$ with $\sk_{\mathrm{root}}$:
\begin{enumerate}[leftmargin=*]
\item $\mathcal{B}_1$ receives a challenge public key $\pk_{\mathrm{root}}$
from its EUF-CMA challenger and signing oracle.
\item $\mathcal{B}_1$ simulates $\mathbf{Exp}^{\mathrm{spoof}}$ for
$\mathcal{A}$, using its EUF-CMA signing oracle to produce
$\sigma_{\mathrm{auth}}$ for each registration. It generates owner
signatures honestly because it knows the per-registration $\sk^\Sig$.
\item When $\mathcal{A}$ outputs $(N^*, \rec^*)$ falling in Case~A,
$\mathcal{B}_1$ outputs $(M^*, \sigma^*_{\mathrm{auth}})$ as its
EUF-CMA forgery.
\end{enumerate}
This succeeds whenever $\mathcal{A}$ wins via Case~A in $G_2$. So
$\Pr[\text{Case A in } G_2] \le \Adv^{\mathrm{euf\text{-}cma}}_{\Sig}(\mathcal{B}_1)$.

\paragraph{Subtlety: owner signature vs.\ authority signature.} The
registration oracle requires that the owner signature
$\sigma_{\mathrm{owner}}$ verify under the freshly generated
$\pk^\Sig$. An adversary who produces a record with a fresh
$\pk^{\Sig,*}$ of its own (not the one in $\mathcal{D}$) and a valid
$\sigma^*_{\mathrm{owner}}$ under it is trivially possible --- the
adversary controls $\pk^{\Sig,*}$. The discriminator is therefore
$\sigma^*_{\mathrm{auth}}$: the authority's counter-signature.
$\mathcal{B}_1$ above targets exactly that signature, which is what
the bound requires. The owner-signature term arises only in the
secondary attack of name re-binding by an evicted owner; we treat that
in the second reduction $\mathcal{B}_2$ below.

\paragraph{Owner re-binding ($\mathcal{B}_2$).} An adversary may
register $N^*$ honestly, learn $\pk^{\Sig,*}$, leak it through a
side-channel, and later claim a re-binding. To prevent counting these
as wins, the spoofing definition requires the record to differ from
$\mathcal{D}[N^*]$ at the \emph{current} time. A re-binding therefore
must produce a fresh $\sigma^*_{\mathrm{auth}}$, again caught by
$\mathcal{B}_1$. The factor $q_R$ in the second term of the theorem
bounds the union over all registered names of the probability that
$\mathcal{A}$ produces a fresh owner signature for one of them,
relevant only when continuity policy demands an old-key signature
during re-registration. The reduction $\mathcal{B}_2$ guesses the
target index $i \in [q_R]$ uniformly and embeds the EUF-CMA challenge
into the $i$-th registration's owner key, paying the factor $q_R$ in
the bound.

\paragraph{Combining.} From $G_0$ to $G_2$ we lose
$p_{\mathrm{root}} + q_H/2^{2\lambda}$. In $G_2$ all wins are Case~A,
which $\mathcal{B}_1$ converts to EUF-CMA forgeries; if continuity is
required we add $q_R \cdot \Adv^{\mathrm{euf\text{-}cma}}_{\Sig}(\mathcal{B}_2)$.
Summing,
\[
\Adv^{\mathrm{spoof}}_{\RNS}(\mathcal{A}) \;\le\;
\Adv^{\mathrm{euf\text{-}cma}}_{\Sig}(\mathcal{B}_1) \;+\;
q_R \cdot \Adv^{\mathrm{euf\text{-}cma}}_{\Sig}(\mathcal{B}_2) \;+\;
p_{\mathrm{root}} \;+\; \frac{q_H}{2^{2\lambda}}. \qedhere
\]
\end{proof}

\begin{remark}
The KEM key-recovery term in the abstract bound is folded into the
case where the adversary somehow possesses $\sk^{\KEM,*}_N$ and uses
it to impersonate. In the formal experiment this never appears
because the win condition is purely record-validity; the KEM
key-recovery term re-emerges in Theorem~2.
\end{remark}

\section{Theorem~2: Connection Binding}

\begin{definition}[Connection-binding experiment]
$\mathbf{Exp}^{\mathrm{cb}}_{\RNS}(\mathcal{A})$:
\begin{enumerate}[leftmargin=*,label=\arabic*.]
\item Setup as in Theorem~1. $\mathcal{A}$ chooses a target name $N^*$,
which the challenger registers honestly via $\mathcal{O}_{\mathrm{reg}}$.
The honest registration produces $(\pk^{\KEM,*}, \sk^{\KEM,*})$; the
challenger keeps $\sk^{\KEM,*}$.
\item An honest client looks up $N^*$, obtains $\pk^{\KEM,*}$, and
runs $\KEM.\Encap(\pk^{\KEM,*}) \to (c^*, K^*)$.
\item The client sends $c^*$ to whatever endpoint resolves $N^*$;
$\mathcal{A}$ controls the network and may interpose, drop, or modify.
\item $\mathcal{A}$ outputs a guess $K'$.
\item Win: $K' = K^*$.
\end{enumerate}
\end{definition}

\begin{theorem}[Connection-binding]\label{thm:cb}
For every polynomial-time $\mathcal{A}$,
\[
\Adv^{\mathrm{cb}}_{\RNS}(\mathcal{A}) \;\le\;
\Adv^{\mathrm{ind\text{-}cca2}}_{\KEM}(\mathcal{B}_3) \;+\;
\Adv^{\mathrm{spoof}}_{\RNS}(\mathcal{B}_4) \;+\; 2^{-\lambda},
\]
where $\mathcal{B}_3, \mathcal{B}_4$ are reductions built from
$\mathcal{A}$ with comparable running time.
\end{theorem}

\begin{proof}
Two cases.

\textbf{Case I: $\mathcal{A}$ does not modify the lookup.} Then the
client uses the honest $\pk^{\KEM,*}$. Recovering $K^*$ is exactly the
IND-CCA2 game: a reduction $\mathcal{B}_3$ embeds its KEM challenge
$(\pk^{\KEM,*}, c^*, \widehat{K})$ ($\widehat{K}$ is the real key in
the IND-CCA2-1 world or random in IND-CCA2-0), runs the experiment, and
outputs 1 if $\mathcal{A}$'s $K'$ equals $\widehat{K}$. This wins
IND-CCA2 with advantage at least $\Pr[\mathcal{A} \text{ wins
Case I}] - 2^{-\lambda}$ (the $-2^{-\lambda}$ accounts for the chance
that $\widehat{K}$ randomly matches $K^*$ in the IND-CCA2-0 world).
$\mathcal{B}_3$ uses its decapsulation oracle to answer any
decapsulation queries the simulation requires, except on $c^*$, which
the experiment forbids.

\textbf{Case II: $\mathcal{A}$ modifies the lookup} so the client
encrypts to a different $\pk'$ controlled by $\mathcal{A}$. Then
$\mathcal{A}$ has produced a record for $N^*$ that the client accepts
as valid but disagrees with the honestly-registered record. By
construction this is a successful spoofing event, so $\mathcal{A}$ has
solved $\mathbf{Exp}^{\mathrm{spoof}}_{\RNS}$. A reduction
$\mathcal{B}_4$ runs $\mathcal{A}$ in this case and outputs the modified
record as its spoofing forgery.

Summing the two cases,
\[
\Adv^{\mathrm{cb}}_{\RNS}(\mathcal{A}) \;\le\;
\Adv^{\mathrm{ind\text{-}cca2}}_{\KEM}(\mathcal{B}_3) \;+\;
\Adv^{\mathrm{spoof}}_{\RNS}(\mathcal{B}_4) \;+\; 2^{-\lambda}. \qedhere
\]
\end{proof}

\begin{corollary}[End-to-end identity guarantee]
Substituting Theorem~1 into Theorem~2,
\begin{align*}
\Adv^{\mathrm{cb}}_{\RNS}(\mathcal{A})
&\le \Adv^{\mathrm{ind\text{-}cca2}}_{\KEM}(\mathcal{B}_3) \\
&\quad + \Adv^{\mathrm{euf\text{-}cma}}_{\Sig}(\mathcal{B}_1)
+ q_R \cdot \Adv^{\mathrm{euf\text{-}cma}}_{\Sig}(\mathcal{B}_2)
+ p_{\mathrm{root}} \\
&\quad + \frac{q_H}{2^{2\lambda}} + 2^{-\lambda}.
\end{align*}
A client encrypting to $N$ reaches only the holder of $\sk_N^\KEM$
unless the KEM is broken, the signature is forged, or the root is
compromised.
\end{corollary}

\section{Theorem~3: Continuity}

\begin{definition}[Continuity experiment]
$\mathbf{Exp}^{\mathrm{cont}}_{\RNS}(\mathcal{A})$: $\mathcal{A}$
honestly registers $N^*$ at time $t_0$, observes the resulting
record. At a later time $t_1$ the challenger's continuity policy
forbids re-registration of $N^*$ without a signature from the prior
$\sk^{\Sig,*}$. $\mathcal{A}$ wins if it produces a fresh record for
$N^*$ accepted by an honest client at time $t_1$ with
$\pk^{\KEM,**} \ne \pk^{\KEM,*}$.
\end{definition}

\begin{theorem}[Continuity]\label{thm:cont}
For every polynomial-time $\mathcal{A}$,
\[
\Adv^{\mathrm{cont}}_{\RNS}(\mathcal{A}) \;\le\;
\Adv^{\mathrm{euf\text{-}cma}}_{\Sig}(\mathcal{B}_5) \;+\;
\Adv^{\mathrm{spoof}}_{\RNS}(\mathcal{B}_6).
\]
\end{theorem}

\begin{proof}
A successful continuity attack is either (a) a forged owner-signature
under $\sk^{\Sig,*}$ (caught by $\mathcal{B}_5$ as an EUF-CMA
forgery), or (b) a record that bypasses the owner check entirely,
which requires the authority to counter-sign without policy
compliance, equivalent to spoofing (caught by $\mathcal{B}_6$).
\end{proof}

\section{Lemma: Sybil Lower Bound}

\begin{definition}[Attestation primitive]
An attestation primitive $\mathcal{T}$ has bypass cost $c(\mathcal{T})
\in \mathbb{R}_{\ge 0}$: the minimum expected resource to produce a
forged attestation accepted by the authority.
\end{definition}

\begin{lemma}[Sybil cost]\label{lem:sybil}
Suppose the authority's policy requires a fresh $\mathcal{T}$-attestation
on every registration. An adversary capable of $n$ Sybil registrations
expends total resources at least $n \cdot c(\mathcal{T}) - O(1)$,
where the $O(1)$ accounts for amortised authority-side costs.
\end{lemma}

\begin{proof}
Each accepted registration carries an attestation that either
(i) was honestly produced under $\mathcal{A}$'s actual identity, in
which case it is accountable and Sybil-free, or (ii) was bypassed,
costing at least $c(\mathcal{T})$. By assumption the adversary controls
no more than one ``honest'' attestation slot, hence at least $n - 1$
attestations are bypassed, totalling $(n-1) \cdot c(\mathcal{T})$
resource. Folding the constant into $O(1)$ yields the claim.
\end{proof}

\begin{corollary}[Sybil-resistance instantiation]
With $\mathcal{T}$ = TPM quote on a measured-boot PCR set, $c(\mathcal{T})$
equals the cost of a hardware root-of-trust compromise (currently
$\gtrsim 10^4$ USD per device per attempt~\cite{rfc7401}); for $n=10^4$
Sybil names the adversary spends at least $10^8$ USD.
\end{corollary}

\section{Concrete Bound for ML-KEM-768 + ML-DSA-65}

Assume $\Adv^{\mathrm{ind\text{-}cca2}}_{\KEM\text{-}768}(\cdot) \le
2^{-128}$ and $\Adv^{\mathrm{euf\text{-}cma}}_{\Sig\text{-}65}(\cdot) \le
2^{-128}$ for adversaries under $2^{128}$ work, $\lambda = 128$.
For an authority handling $q_R = 2^{30}$ registrations and
$q_H = 2^{60}$ hash queries:
\begin{align*}
\Adv^{\mathrm{spoof}}_{\RNS}(\mathcal{A})
&\le 2^{-128} + 2^{30} \cdot 2^{-128} + p_{\mathrm{root}} + 2^{60-256} \\
&\le 2^{-97} + p_{\mathrm{root}}
\end{align*}
\begin{align*}
\Adv^{\mathrm{cb}}_{\RNS}(\mathcal{A})
&\le 2^{-128} + 2^{-97} + p_{\mathrm{root}} + 2^{-128} \\
&\le 2^{-97} + p_{\mathrm{root}}
\end{align*}
With $p_{\mathrm{root}} = 2^{-40}$ (a generous root-key-compromise
budget) the dominant term is the policy-controlled
$p_{\mathrm{root}}$, not any cryptographic primitive. \RNS{}'s
security is therefore root-policy-bound, exactly as designed.

\section{Discussion}

\paragraph{Tightness.} The factor $q_R$ in Theorem~1's second term is
the standard guess-the-target loss in EUF-CMA reductions. It can be
removed using a multi-user EUF-CMA assumption~\cite{boneh2001ibe},
which is satisfied by ML-DSA in the random-oracle model. Substituting
multi-user EUF-CMA the bound becomes
$\Adv^{\mathrm{spoof}}_{\RNS}(\mathcal{A}) \le 2 \cdot
\Adv^{\mathrm{m\text{-}euf\text{-}cma}}_{\Sig}(\mathcal{B}) +
p_{\mathrm{root}} + q_H/2^{2\lambda}$.

\paragraph{Random-oracle reliance.} The collision step of $G_2$ uses
the random-oracle property of $H$. Substituting a collision-resistant
hash (SHA3-256) the same bound follows from the collision-resistance
advantage instead of the birthday term, with no qualitative change.

\paragraph{Forward secrecy.} Long-term KEM keys mean a compromise of
$\sk^\KEM$ at time $t > t_0$ exposes all prior session keys. The
remedy is ephemeral KEM encapsulation atop the long-term identity:
the long-term key authenticates the ephemeral, the ephemeral provides
forward secrecy. This is the HPKE pattern~\cite{rfc9180} and is
orthogonal to the binding analysed here.

\paragraph{Post-quantum.} The proofs above treat $\KEM$ and $\Sig$ as
abstract primitives. Substituting ML-KEM-768 and ML-DSA-65 yields a
post-quantum scheme whose security holds against quantum adversaries
modulo standard lattice assumptions. The random-oracle step is
quantum-accessible (QROM); the same bound holds with the standard
$O(\sqrt{q_H})$ factor.

\paragraph{Auditability.} The single root key is the only unconditional
trust assumption. Operators can split $\sk_{\mathrm{root}}$ into a
threshold via an MPC scheme, log every $\sigma_{\mathrm{auth}}$ to a
transparency log~\cite{certtransparency}, and gate root signing on
human approval. Each of these reduces $p_{\mathrm{root}}$ further; the
bounds above hold for any $p_{\mathrm{root}}$.

\bibliographystyle{plain}
\begin{thebibliography}{9}
\bibitem{shoup2001}
V.~Shoup. ``OAEP Reconsidered.'' \emph{CRYPTO}, 2001.
\bibitem{boneh2001ibe}
D.~Boneh and M.~Franklin. ``Identity-Based Encryption from the Weil
Pairing.'' \emph{CRYPTO}, 2001.
\bibitem{rfc7401}
R.~Moskowitz et al. ``Host Identity Protocol Version 2 (HIPv2).''
RFC~7401, 2015.
\bibitem{rfc9180}
R.~Barnes et al. ``Hybrid Public Key Encryption.'' RFC~9180, 2022.
\bibitem{certtransparency}
B.~Laurie, A.~Langley, E.~Kasper. ``Certificate Transparency.''
RFC~6962, 2013.
\end{thebibliography}

\end{document}
